% Part: methods
% Chapter: proofs
% Section: resources

\documentclass[../../../include/open-logic-section]{subfiles}

\begin{document}

\olfileid[pt]{mth}{prf}{res}

\section{Outros Recursos}

Há muitos livros sobre como fazer provas em matemática que podem ser
úteis. Confira em particular \emph{How to Read and do Proofs: An
Introduction to Mathematical Thought Processes} \citep{Solow2013} e
\emph{How to Prove It: A Structured Approach} \citep{Velleman2019}.  O
\href{http://www.people.vcu.edu/~rhammack/BookOfProof/BookOfProof.pdf}{\emph{Book
of Proof}} \citep{Hammack2013} e o
\href{https://scholarworks.gvsu.edu/books/7/}{\emph{Mathematical
Reasoning}} \citep{Sandstrum2019} são livros sobre provas que estão
disponíveis gratuitamente online.  Filósofos podem achar
\emph{More
Precisely: The Math you need to do Philosophy} \citep{Steinhart2018}
uma boa introdução ao raciocínio matemático.

Há também vários guias mais curtos sobre provas disponíveis na
internet; por exemplo,
\href{https://math.berkeley.edu/~hutching/teach/proofs.pdf}{``Introduction
  to Mathematical Arguments''} \citep{Hutchings2003} e
\href{https://eugeniacheng.com/wp-content/uploads/2017/02/cheng-proofguide.pdf}{``How to
  write proofs''} \citep{Cheng2004}.

\subsection{Vídeos Motivacionais}

Sente que não tem motivação para fazer sua tarefa de casa? Desanimado?
Estes vídeos podem ajudar!

\begin{itemize}
\item \url{https://www.youtube.com/watch?v=ZXsQAXx_ao0}
\item \url{https://www.youtube.com/watch?v=BQ4yd2W50No}
\item \url{https://www.youtube.com/watch?v=StTqXEQ2l-Y}
\end{itemize}

\end{document}
